\chapter{Liczby porządkowe i rekurencja pozaskończona}\label{ch:ordinals}
\chaptermark{Liczby porządkowe i rekurencja pozaskończona}

\section{Od dobrych porządków do obiektów kanonicznych}
Dwa dobrze uporządkowane zbiory mogą być izomorficzne, mimo że mają różne elementy. Chcemy reprezentować ich typy porządkowe konkretnymi zbiorami. W tym rozdziale pracujemy w ZF, o ile nie wskazano inaczej; przy konstrukcji liczb kardynalnych jako najmniejszych ordynałów danej mocy korzystamy z AC.
\begin{definicja}[Zbiór przechodni i liczba porządkowa]\label{def:ordinal}
Zbiór $X$ jest \emph{przechodni}, jeżeli $y\in x\in X$ implikuje $y\in X$. \emph{Liczbą porządkową} (ordynałem) nazywamy zbiór przechodni $\alpha$, na którym relacja $\in$ jest ścisłym dobrym porządkiem: jest przechodnia, dowolne dwa różne elementy są porównywalne i każdy niepusty podzbiór ma element najmniejszy względem $\in$.
\end{definicja}
\begin{przyklad}[Pierwsze ordynały]
$0=\varnothing$, $1=\{0\}$, $2=\{0,1\}$, $3=\{0,1,2\}$ są ordynałami. Również $\omega$, skonstruowane w Twierdzeniu~\ref{thm:zf-omega},
jest ordynałem: jego elementy to skończone ordynały von Neumanna.
Lemat~\ref{lem:zf-vn-successor} dowodzi przechodniości $\omega$,
a Twierdzenie~\ref{thm:zf-wellordered-omega} wraz
z Lematem~\ref{lem:zf-initial-segments} pokazuje,
że relacja $\in$ dobrze porządkuje $\omega$. Ordynał $\omega+1=\omega\cup\{\omega\}$ ma element największy, podczas gdy $\omega$ go nie ma. Dlatego te dwa porządki nie są izomorficzne, choć ich zbiory są równoliczne.
\end{przyklad}
\begin{lemat}[Elementy ordynału i odcinki początkowe]\label{lem:ordinal-segments}
Jeżeli $\alpha$ jest ordynałem i $\beta\in\alpha$, to $\beta$ jest ordynałem oraz $\beta=\{\gamma\in\alpha:\gamma\in\beta\}$. Każdy właściwy odcinek początkowy $\alpha$ ma postać $\beta$ dla dokładnie jednego $\beta\in\alpha$.
\end{lemat}
\begin{proof}
Ponieważ $\alpha$ jest przechodni, $\beta\subseteq\alpha$. Dla $\gamma\in\beta$ i $\delta\in\gamma$ przechodniość $\in$ na $\alpha$ daje $\delta\in\beta$; zatem $\beta$ jest przechodni. Ograniczenie dobrego porządku $\in$ z $\alpha$ do $\beta$ pozostaje dobrym porządkiem, więc $\beta$ jest ordynałem. Równość zbiorów wynika z $\beta\subseteq\alpha$. Jeśli $I$ jest właściwym odcinkiem początkowym, zbiór $\alpha\setminus I$ ma najmniejszy element $\beta$; wtedy $I=\{\gamma\in\alpha:\gamma\in\beta\}=\beta$. Jednoznaczność wynika z antysymetrii porządku liniowego.
\end{proof}
\begin{twierdzenie}[Porównanie ordynałów]\label{thm:ordinal-trichotomy}
Dla dowolnych ordynałów $\alpha,\beta$ zachodzi dokładnie jedna z relacji $\alpha\in\beta$, $\alpha=\beta$, $\beta\in\alpha$.
\end{twierdzenie}
\begin{proof}
Lemat~\ref{lem:choice-wellorder-comparison} porównuje dobre porządki $(\alpha,\in)$ i $(\beta,\in)$. Pokażemy najpierw, że izomorfizm $f$ ordynału $\alpha$ na odcinek początkowy ordynału $\beta$ musi być identycznością. Gdyby nie, weźmy najmniejsze $\gamma\in\alpha$ z $f(\gamma)\ne\gamma$. Dla każdego $\delta\in\gamma$ mamy $f(\delta)=\delta$, a izomorfizm przenosi zbiór poprzedników $\gamma$ na zbiór poprzedników $f(\gamma)$. Z Lematu~\ref{lem:ordinal-segments} otrzymujemy $\gamma=f(\gamma)$, sprzeczność. Izomorficzne ordynały są więc równe, a izomorfizm z właściwym odcinkiem daje odpowiednio $\alpha\in\beta$ lub $\beta\in\alpha$. Rozłączność wynika z asymetrii $\in$ na ordynałach i aksjomatu ufundowania.
\end{proof}
\begin{wniosek}[Następnik i supremum ordynałów]\label{cor:ordinal-sup}
Dla ordynału $\alpha$ zbiór $\alpha+1:=\alpha\cup\{\alpha\}$ jest ordynałem i jest najmniejszym ordynałem większym od $\alpha$. Dla każdego zbioru $S$ ordynałów jego suma $\sup S:=\bigcup S$ jest ordynałem i najmniejszym wspólnym ograniczeniem górnym w porządku $\le$, gdzie $\alpha\le\beta$ oznacza $\alpha\subseteq\beta$.
\end{wniosek}
\begin{proof}
\emph{Krok 1: zagnieżdżenie i suma.}
Lemat~\ref{lem:ordinal-segments}
i Twierdzenie~\ref{thm:ordinal-trichotomy}
dają dla dowolnych ordynałów
$\alpha,\beta$ zawieranie
$\alpha\subseteq\beta$ lub
$\beta\subseteq\alpha$.
Połóżmy $\gamma:=\bigcup S$.
Jeśli $x\in y\in\gamma$, to
dla pewnego $\delta\in S$
mamy $y\in\delta$; stąd
$x\in\delta\subseteq\gamma$.
Zatem $\gamma$ jest przechodni.
Dla różnych $x,y\in\gamma$
istnieją $\delta,\varepsilon\in S$
z $x\in\delta$, $y\in\varepsilon$.
Zagnieżdżenie tych dwóch ordynałów
daje wspólny ordynał zawierający
$x,y$, gdzie porównuje je relacja $\in$.
Relacja $\in$ na $\gamma$
jest więc ścisłym porządkiem liniowym.

\emph{Krok 2: dobry porządek.}
Jeśli $X\subseteq\gamma$ jest niepusty,
ustalmy $x\in X$ i $\delta\in S$
z $x\in\delta$. Niepusty zbiór
$X\cap\delta$ ma najmniejszy
element $m$, gdyż $\delta$ jest
dobrze uporządkowany.
Jeżeli $y\in X$ i $y\in m$,
to $m\in\delta$ i przechodniość
$\delta$ dają $y\in X\cap\delta$,
wbrew minimalności $m$.
Zatem $m$ jest najmniejszym
elementem całego $X$.
Dla $S=\varnothing$ mamy
$\gamma=\varnothing$, również ordynał.

\emph{Krok 3: supremum.}
Z definicji sumy zbiorów
każde $\delta\in S$ zawiera się
w $\gamma$, a każdy ordynał
zawierający wszystkie $\delta\in S$
zawiera także $\gamma$.
To dowodzi $\gamma=\sup S$.

\emph{Krok 4: następnik.}
Zbiór $\alpha+1:=\alpha\cup\{\alpha\}$
jest przechodni: jeśli
$y\in x\in\alpha+1$, to
albo $x\in\alpha$ i $y\in\alpha$
przez przechodniość $\alpha$,
albo $x=\alpha$ i także
$y\in\alpha$.
Relacja $\in$ porządkuje liniowo
$\alpha+1$: elementy $\alpha$
są porównywalne i wszystkie
poprzedzają $\alpha$.
Niepusty podzbiór przecinający
$\alpha$ ma minimum w $\alpha$,
a jedyny niepusty podzbiór
nieprzecinający $\alpha$ to
$\{\alpha\}$.
Zatem $\alpha+1$ jest ordynałem.
Jeśli $\alpha<\eta$ dla ordynału
$\eta$, to $\alpha\in\eta$,
a z przechodniości $\eta$
wynika $\alpha+1\subseteq\eta$.
Nie ma więc ordynału
ściśle pomiędzy $\alpha$ i $\alpha+1$.
\end{proof}
\begin{definicja}[Ordynały graniczne]\label{def:ordinal-limit}
Ordynał $\lambda$ jest \emph{graniczny}, jeśli $\lambda\ne0$ i nie jest następnikiem żadnego ordynału. Wtedy $\lambda=\sup\{\beta:\beta<\lambda\}$. Przykładem jest $\omega$, podczas gdy $\omega+1$ jest następnikiem.
\end{definicja}


\begin{lemat}[Ordynały graniczne jako sumy poprzedników]
\label{lem:ordinal-limit-union}
Jeżeli $\lambda$ jest niezerowym ordynałem granicznym, to
$\lambda=\bigcup_{\beta<\lambda}\beta$.
Jeżeli $\lambda=\gamma+1$ jest następnikiem, to
$\bigcup_{\beta<\lambda}\beta=\gamma$.
\end{lemat}
\begin{proof}
Każdy $\beta\in\lambda$ jest podzbiorem $\lambda$
na mocy Lematu~\ref{lem:ordinal-segments},
zatem $U:=\bigcup_{\beta<\lambda}\beta\subseteq\lambda$.
Załóżmy, że $\lambda$ jest graniczny i
istnieje $x\in\lambda\setminus U$.
Gdyby istniał $y\in\lambda$ z $x\in y$,
to $x\in y\subseteq U$, sprzeczność.
Zatem $x$ jest największym elementem
dobrze uporządkowanego $\lambda$
(właściwa liniowość relacji $\in$ wynika
z definicji ordynału).
Każdy $y\in\lambda$ jest więc równy $x$
lub należy do $x$, skąd
$\lambda=x\cup\{x\}=x+1$,
wbrew założeniu graniczności.
Stąd $\lambda=U$.

Jeżeli $\lambda=\gamma+1$, każdy $\beta\in\lambda$
należy do $\gamma$ lub jest równy $\gamma$,
a więc $\beta\subseteq\gamma$.
Z drugiej strony $\gamma\in\lambda$,
więc $U=\gamma$.
\end{proof}
\begin{przyklad}[Pusta suma i pierwszy ordynał graniczny]
\label{ex:ordinal-suprema}
Dla pustego zbioru ordynałów
$\sup\varnothing=\bigcup\varnothing=0$.
Dla $S=\{1,3\}$ otrzymujemy $\sup S=3$,
ponieważ $1\subseteq3$ i $1\cup3=3$.
Ponadto $\sup\{n:n<\omega\}=\omega$:
dla każdego $m\in\omega$ istnieje
następnik $m+1\in\omega$ zawierający $m$.
Natomiast
$\sup\{\beta:\beta<\omega+1\}=\omega$,
a nie $\omega+1$, zgodnie z lematem.
\end{przyklad}

\section{Indukcja pozaskończona}
\begin{twierdzenie}[Zasada indukcji pozaskończonej]\label{thm:ordinal-induction}
Niech $\alpha$ będzie ordynałem, a $P(\beta)$
formułą języka ZF z ustalonymi parametrami
i zmienną wolną $\beta$, rozważaną dla $\beta<\alpha$. Jeśli dla każdego $\beta<\alpha$ z założenia $P(\gamma)$ dla wszystkich $\gamma<\beta$ wynika $P(\beta)$, to $P(\beta)$ zachodzi dla wszystkich $\beta<\alpha$.
\end{twierdzenie}
\begin{proof}
Gdyby istniał kontrprzykład, zbiór
$E=\{\beta\in\alpha:\neg P(\beta)\}$
istniałby przez separację i byłby niepusty. Niech $\beta$ będzie jego najmniejszym elementem. Wszystkie $\gamma<\beta$ spełniają $P$, a założenie kroku daje $P(\beta)$, sprzeczność.
\end{proof}
\begin{uwaga}[Trzy przypadki w praktyce]
Wygodnie rozdziela się dowód na przypadek $0$, następnik $\beta+1$ i ordynał graniczny $\lambda$. Nie są to trzy niezależne aksjomaty: stanowią sposób sprawdzenia jednego warunku indukcji z poprzedniego twierdzenia.
\end{uwaga}

\section{Rekurencja pozaskończona}
\begin{twierdzenie}[Rekurencja na ustalonym ordynale]\label{thm:ordinal-recursion}
Niech $\alpha$ będzie ordynałem, a $F$ operacją klasową zdefiniowaną formułą teorii mnogości z ustalonymi parametrami, która każdej funkcji $g$ o dziedzinie będącej ordynałem $\beta<\alpha$ przypisuje dokładnie jeden zbiór $F(g)$. Wówczas istnieje dokładnie jedna funkcja $h$ o dziedzinie $\alpha$, dla której
\[
\forall\beta<\alpha\quad h(\beta)=F(h|_\beta).
\]
W dowodzie korzystamy ze schematu zastępowania; nie potrzeba aksjomatu wyboru.
\end{twierdzenie}
\begin{proof}
\emph{Krok 1: przybliżenia i ich zgodność.}
Poprawnym przybliżeniem długości
$\beta\le\alpha$ nazywamy funkcję
$g$ o dziedzinie $\beta$ taką, że
$g(\gamma)=F(g|_\gamma)$ dla każdego $\gamma<\beta$.
Jest to własność wyrażalna formułą
języka ZF z ustalonymi parametrami.
Jeśli dwa przybliżenia $g,k$
różniłyby się na wspólnej dziedzinie,
istniałby najmniejszy punkt $\gamma$ różnicy.
Wtedy $g|_\gamma=k|_\gamma$, więc
\[
g(\gamma)=F(g|_\gamma)=F(k|_\gamma)=k(\gamma),
\]
sprzeczność.
Dla danej dziedziny istnieje co najwyżej
jedno przybliżenie.

\emph{Krok 2: indukcyjne istnienie.}
Stosujemy indukcję pozaskończoną
do ordynału $\alpha+1$, aby uwzględnić
wszystkie $\beta\le\alpha$.
Dla $\beta=0$ funkcja pusta jest przybliżeniem.
Jeśli $\beta=\delta+1\le\alpha$,
to z hipotezy indukcyjnej mamy
jednoznaczne przybliżenie $g$
o dziedzinie $\delta$.
Ponieważ $\delta<\alpha$,
operacja $F(g)$ jest określona
i jest zbiorem.
Funkcja $g\cup\{(\delta,F(g))\}$
jest poprawnym przybliżeniem długości $\beta$.

Niech $\beta\le\alpha$ będzie
niezerowym ordynałem granicznym.
Dla każdego $\delta<\beta$
istnieje dokładnie jedno
przybliżenie $g_\delta$ długości $\delta$.
Ponieważ jest to przyporządkowanie
jednoznacznie definiowalne formułą,
schemat zastępowania daje zbiór
$\mathcal G=\{g_\delta:\delta<\beta\}$.
Poprzedni krok pokazuje zgodność
przybliżeń, więc $g:=\bigcup\mathcal G$
jest funkcją.
Na mocy Lematu~\ref{lem:ordinal-limit-union}
jej dziedzina jest równa
$\bigcup_{\delta<\beta}\delta=\beta$.
Dla $\gamma<\beta$ mamy
$\gamma+1<\beta$ (gdyż $\beta$ jest graniczny),
zatem $g|_{\gamma+1}=g_{\gamma+1}$.
Poprawność tego przybliżenia
daje $g(\gamma)=F(g|_\gamma)$.
Wobec tego $g$ jest poprawne
na całej dziedzinie $\beta$.

\emph{Krok 3: zakończenie.}
Indukcja daje poprawne przybliżenie
$h$ długości $\alpha$.
Spełnia ono wymagane równanie,
a zgodność dowolnych poprawnych
przybliżeń daje jego jednoznaczność.
Ani w etapie następnikowym, ani granicznym
nie dokonuje się nieskończonego
wyboru spośród wielu dopuszczalnych wartości.
\end{proof}
\begin{przyklad}[Dodawanie ordynałów]\label{ex:ordinal-addition}
Dla ustalonego ordynału $\alpha$ definiujemy
\emph{sumę porządkową} rekurencyjnie
po drugim argumencie:
\[
\alpha+0:=\alpha,\qquad
\alpha+(\beta+1):=(\alpha+\beta)+1,
\]
a dla niezerowego ordynału granicznego $\lambda$:
\[
\alpha+\lambda:=\sup_{\beta<\lambda}(\alpha+\beta).
\]
Symbol $+1$ po prawej oznacza następnik ordynału
z Wniosku~\ref{cor:ordinal-sup}.
Są to trzy przypadki jednoznacznej operacji
na dotychczas wyliczonych wartościach.
Dla dowolnego ustalonego ordynału $\theta$
Twierdzenie~\ref{thm:ordinal-recursion}
konstruuje funkcję $\beta\mapsto\alpha+\beta$
na $\beta<\theta$:
operację klasową na \emph{wszystkich}
funkcjach o dziedzinie $\beta$
można zdefiniować tą samą regułą,
a w przypadku niewłaściwych wartości
zwracać umownie $0$.
Indukcja zapewnia, że na konstruowanej funkcji
wszystkie wartości są ordynałami.
Jeśli $\theta<\theta'$, obie konstrukcje
zgadzają się na wspólnej dziedzinie
przez jednoznaczność rekurencji.

Dla każdego $n\in\omega$ indukcja
daje $1+n=n+1$:
dla $n=0$ mamy $1+0=1$,
a z $1+n=n+1$ otrzymujemy
$1+(n+1)=(n+1)+1$.
Wobec tego
\[
1+\omega=\sup_{n<\omega}(1+n)
=\sup_{n<\omega}(n+1)=\omega.
\]
Natomiast $\omega+1>\omega$
z definicji następnika.
Dodawanie ordynałów nie jest więc przemienne.
\end{przyklad}


\begin{przyklad}[Hierarchia von Neumanna]
\label{ex:ordinal-vn-hierarchy}
Niech $\alpha$ będzie ordynałem z $4\le\alpha$.
Definiujemy rekurencyjnie $V_\beta$ dla $\beta<\alpha$:
$V_0:=\varnothing$,
$V_{\delta+1}:=\mathcal P(V_\delta)$
wtedy, gdy $\delta+1<\alpha$,
oraz $V_\lambda:=\bigcup_{\beta<\lambda}V_\beta$
dla niezerowych ordynałów granicznych $\lambda<\alpha$.
Warunki na indeksy gwarantują, że
każdy symbol po lewej stronie oznacza
poziom należący do rozważanej rodziny.
Twierdzenie~\ref{thm:ordinal-recursion}
gwarantuje istnienie funkcji
$\beta\mapsto V_\beta$ na ustalonym zbiorze
indeksów $\alpha$.
W każdym kroku otrzymujemy jednoznacznie zbiór,
więc nie używamy aksjomatu wyboru.
W szczególności
$V_1=\{\varnothing\}$,
$V_2=\{\varnothing,\{\varnothing\}\}$,
a $V_3=\mathcal P(V_2)$ ma cztery elementy.
W ZF nie zakładamy istnienia jednego zbioru
zawierającego poziomy $V_\beta$ dla
\emph{wszystkich} ordynałów $\beta$.
\end{przyklad}

\section{Mnożenie ordynałów i różnica względem arytmetyki mocy}

\begin{definicja}[Iloczyn porządkowy ordynałów]
\label{def:ordinal-multiplication}
Dla dowolnych ordynałów $\alpha,\beta$
iloczyn porządkowy $\alpha\cdot\beta$
definiujemy rekurencyjnie po $\beta$:
\[
\alpha\cdot0:=0,\qquad
\alpha\cdot(\beta+1):=(\alpha\cdot\beta)+\alpha,
\]
a dla niezerowego ordynału granicznego $\lambda$:
\[
\alpha\cdot\lambda:=\sup_{\beta<\lambda}
(\alpha\cdot\beta).
\]
Wzór zawiera dodawanie \emph{ordynałów}
z Przykładu~\ref{ex:ordinal-addition},
a nie dodawanie mocy.
Na każdym ustalonym odcinku ordynałów
istnienie i jednoznaczność funkcji
$\beta\mapsto\alpha\cdot\beta$
wynikają z Twierdzenia~\ref{thm:ordinal-recursion}.
Aby zastosować twierdzenie dokładnie w jego postaci,
dla niepoprawnych przybliżeń (np. z wartościami
niebędącymi ordynałami) definiujemy wartość
pomocniczej operacji klasowej jako $0$.
Na właściwych przybliżeniach rekurencja
stosuje wypisane powyżej równania
i indukcja zachowuje przynależność do ordynałów.
Dzięki jednoznaczności wynik nie zależy
od wybranego odcinka indeksów.
\end{definicja}
\begin{przyklad}[Nieprzemienność mnożenia ordynałów]
\label{ex:ordinal-mul-noncommutative}
Dla każdego $n\in\omega$ rekurencja
daje $2\cdot n=2n$ w arytmetyce
skończonych ordynałów, zatem
$2\cdot\omega=\sup_{n<\omega}2n=\omega$.
Natomiast
$\omega\cdot2=(\omega\cdot1)+\omega
=\omega+\omega>\omega$.
Istotnie, $\omega+1$ należy do rodziny
$\{\omega+n:n<\omega\}$,
której supremum wynosi $\omega+\omega$,
więc $\omega+\omega\ge\omega+1>\omega$.
Odpowiada to ustawieniu drugiej kopii
porządku $\omega$ po wszystkich punktach pierwszej.
Wobec tego $2\cdot\omega\ne\omega\cdot2$,
chociaż oba ordynały mają moc $\aleph_0$.
W arytmetyce mocy
$2\cdot\aleph_0=\aleph_0\cdot2=\aleph_0$
(Twierdzenie~\ref{thm:card-n2}).
\end{przyklad}

\section{Liczby kardynalne jako ordynały początkowe}
\begin{definicja}[Ordynał początkowy]\label{def:initial-ordinal}
Ordynał $\kappa$ jest \emph{początkowy}, jeżeli żaden ordynał $\beta<\kappa$ nie jest z nim równoliczny. W ZFC \emph{liczbą kardynalną} zbioru $A$ nazywamy najmniejszy ordynał równoliczny z $A$ i oznaczamy $|A|$.
\end{definicja}
\begin{twierdzenie}[Istnienie reprezentanta kardynalnego]\label{thm:initial-exists}
Przy założeniu ZFC każdy zbiór $A$ jest równoliczny z dokładnie jednym ordynałem początkowym.
\end{twierdzenie}
\begin{proof}
\emph{Krok 1: dobry porządek na $A$.}
Na mocy Twierdzenia~\ref{thm:choice-wellorder}
istnieje dobry porządek $\prec$ na $A$.
Jest to jedyne miejsce w konstrukcji,
w którym wykorzystujemy AC.
Dalszy argument działa dla dowolnego ustalonego
zbioru dobrze uporządkowanego $(W,\prec)$ w ZF.

\emph{Krok 2: rekursja wzdłuż ustalonego dobrego porządku.}
Dla $w\in W$ oznaczmy
$W_{<w}:=\{v\in W:v\prec w\}$ oraz
$W_{\le w}:=W_{<w}\cup\{w\}$.
Chcemy skonstruować funkcję $r$ na $W$
spełniającą
\[
r(w)=\{r(v):v\prec w\}.
\tag{$*$}
\]
Poprawnym $w$-przybliżeniem nazwijmy funkcję
$f$ o dziedzinie $W_{\le w}$,
która dla każdego $t\preceq w$
(tj. $t\prec w$ albo $t=w$) spełnia
$f(t)=\{f(v):v\prec t\}$.
Nie zakładamy istnienia wspólnego z góry
zbioru $V$ dla wszystkich przybliżeń:
warunek jest formułą ZF o wykresie funkcji.

Dwa poprawne przybliżenia są zgodne
na wspólnej dziedzinie.
Gdyby się różniły, zbiór punktów niezgodności
miałby najmniejszy element $t$.
Wartości przy wszystkich $v\prec t$
byłyby równe, toteż równanie ($*$)
narzucałoby tę samą wartość w $t$,
co przeczy wyborowi $t$.
Na ustalonej dziedzinie istnieje więc
co najwyżej jedno poprawne przybliżenie.

Pokażemy indukcją po dobrym porządku $W$,
że dla każdego $w\in W$ istnieje
poprawne $w$-przybliżenie $f_w$.
Załóżmy, że istnieje ono dla wszystkich $v\prec w$.
Dzięki jednoznaczności schemat zastępowania
daje zbiór
$\mathcal F_w:=\{f_v:v\prec w\}$.
Suma $u:=\bigcup\mathcal F_w$ jest funkcją:
wszystkie składniki są zgodne.
Jej dziedziną jest dokładnie $W_{<w}$,
ponieważ każdy $t\prec w$ należy do
dziedziny $f_t$, a żadna dziedzina
$f_v$ nie zawiera elementów spoza $W_{<w}$.
Dla każdego $t\prec w$ funkcja $u$
zgadza się z $f_t$ na $W_{\le t}$,
zatem spełnia równanie ($*$) w punkcie $t$.
Kładziemy
\[
f_w:=u\cup\{(w,u[W_{<w}])\}.
\]
Nowy argument $w$ nie należał do
dziedziny $u$, więc jest to funkcja
na $W_{\le w}$ spełniająca ($*$).
Dla najmniejszego $w$ mamy
$\mathcal F_w=\varnothing$ i $u=\varnothing$,
a więc konstrukcja obejmuje przypadek początkowy.

Indukcja jest zasadna bez AC:
jeżeli istniałby kontrprzykład w $W$,
miałby element najmniejszy, dla którego
wszystkie wcześniejsze przybliżenia
już istnieją, co daje sprzeczność.
Zastępowanie zbiera rodzinę
$\{f_w:w\in W\}$ w zbiór.
Jej suma $r:=\bigcup_{w\in W}f_w$
jest funkcją o dziedzinie $W$
spełniającą ($*$).
Gdy $W=\varnothing$, bierzemy funkcję pustą.

\emph{Krok 3: obraz funkcji jest ordynałem.}
Dla $u\prec v$ zachodzi $r(u)\in r(v)$
na mocy ($*$). Jeśli $r(u)=r(v)$
dla różnych $u,v\in W$, to jeden
z nich poprzedza drugi i otrzymujemy
$r(u)\in r(u)$, wbrew aksjomatowi ufundowania.
Zatem $r$ jest iniekcją.
W drugą stronę, jeśli $r(u)\in r(v)$,
to na mocy ($*$) istnieje $t\prec v$
z $r(u)=r(t)$. Iniektywność daje
$u=t\prec v$.
Wykazaliśmy więc
\[
u\prec v\quad\Longleftrightarrow\quad
r(u)\in r(v).
\tag{$**$}
\]

Dla $w\in W$ zbiór
$r(w)=\{r(v):v\prec w\}$ jest przechodni:
jeżeli $x\in r(v)\in r(w)$,
to dla pewnego $t\prec v\prec w$
mamy $x=r(t)\in r(w)$.
Na mocy ($**$) relacja $\in$
na $r(w)$ jest dobrym porządkiem,
przeniesionym z $W_{<w}$ przez iniekcję $r$.
Zatem każde $r(w)$ jest ordynałem.
Zbiór $\gamma:=r[W]$ istnieje
przez zastępowanie, jest przechodni
z równania ($*$) i jest dobrze
uporządkowany przez $\in$
na mocy ($**$).
Jest więc ordynałem, a
$r:W\to\gamma$ izomorfizmem
porządkowym.
Jednoznaczność $\gamma$ wynika
z Twierdzenia~\ref{thm:ordinal-trichotomy}:
jedyny izomorfizm dwóch ordynałów
jest identycznością.

\emph{Krok 4: najmniejszy reprezentant mocy.}
Wracając do $W=A$, mamy ordynał
$\alpha$ równoliczny z $A$.
Zbiór
\[
E:=\{\beta\in\alpha+1:\beta\approx A\}
\]
istnieje przez separację i jest niepusty,
ponieważ $\alpha\in E$.
Dobry porządek na $\alpha+1$
daje najmniejszy element $\kappa$ zbioru $E$.
Jest on początkowy: gdyby
$\beta<\kappa$ był równoliczny z $\kappa$,
byłby również równoliczny z $A$,
a zatem $\beta\in E$, wbrew
minimalności $\kappa$.

Jeżeli $\kappa'$ byłby innym ordynałem
początkowym równolicznym z $A$,
to Twierdzenie~\ref{thm:ordinal-trichotomy}
dawałoby $\kappa<\kappa'$ lub $\kappa'<\kappa$.
Mniejszy byłby równoliczny z większym,
co przeczyłoby początkowości większego.
Stąd reprezentant kardynalny
istnieje i jest jednoznaczny.
\end{proof}
\begin{uwaga}[Dwie różne miary nieskończoności]
Ordynał mierzy \emph{typ dobrego porządku}, liczba kardynalna — \emph{liczbę elementów}. Zbiory $\omega$ i $\omega+1$ są równoliczne:
bijekcję $\omega\to\omega+1$ wyznaczają
wzory $f(0)=\omega$ oraz $f(n+1)=n$.
Nie są jednak izomorficzne jako dobre porządki,
gdyż $\omega$ nie ma elementu największego,
a $\omega+1$ ma.
Najmniejszą nieskończoną liczbą kardynalną
jest $\omega$, oznaczany również $\aleph_0$:
każdy ordynał mniejszy od $\omega$
jest skończony, a $\omega$ jest nieskończony,
bo każdy jego element ma większy następnik.
\end{uwaga}

\section{Dlaczego zbiór wszystkich ordynałów nie istnieje?}
\begin{twierdzenie}[Paradoks Burali-Forti]\label{thm:burali-forti}
Nie istnieje zbiór $O$, którego elementami są dokładnie wszystkie ordynały.
\end{twierdzenie}
\begin{proof}
Przypuśćmy, że taki zbiór $O$ istnieje. Suma $\gamma=\bigcup O$ jest ordynałem na mocy Wniosku~\ref{cor:ordinal-sup}. Z definicji $O$ wynika $\gamma\in O$, zatem $\gamma\subseteq\bigcup O=\gamma$, co jeszcze nie daje sprzeczności. Ale $\gamma+1$ również jest ordynałem, więc $\gamma+1\in O$ i $\gamma+1\subseteq\gamma$. Tymczasem $\gamma\in\gamma+1$, skąd $\gamma\in\gamma$, wbrew aksjomatowi ufundowania. Zatem $O$ nie jest zbiorem.
\end{proof}
\begin{uwaga}
Wyrażenie „klasa wszystkich ordynałów” oznacza tu własność definiowalną formułą języka ZF, a nie dodatkowy obiekt, po którym wolno kwantyfikować w języku pierwszego rzędu.
\end{uwaga}